\chapter{Experiments and Evaluation Protocol}
\label{ch:experiments}

The current repository does \emph{not} ship a fixed set of
end-to-end benchmark experiments with recorded numerical results.
The migration to a Streamlit / local-first architecture prioritised
correctness and reproducibility over benchmark curation. This chapter
therefore defines a rigorous experimental \emph{protocol} that a
future user of the platform can execute to produce the numbers.
Placeholders are labelled \exper{} throughout; when the protocol is
executed, the numbers can be written into Appendix~\ref{app:extra}
without altering the rest of the report.

\section{Datasets}
\label{sec:datasets}

The recommended benchmark set is intentionally small and
tabular-only, chosen to be freely available, well studied, and
diverse in task type and size.

\begin{table}[H]
\centering
\caption{Recommended benchmark datasets. All are freely available
from the UCI Machine Learning Repository \citep{uci}.}
\label{tab:datasets}
\begin{tabular}{@{}lrrll@{}}
\toprule
\textbf{Dataset} & \textbf{Rows} & \textbf{Cols} & \textbf{Target} & \textbf{Task} \\
\midrule
Iris \citep{uciiris}       & 150   & 4  & Species       & Multi-class classification \\
Wine \citep{uciwine}       & 178   & 13 & Class         & Multi-class classification \\
Breast Cancer \citep{ucibreast} & 569 & 30 & Diagnosis     & Binary classification \\
Adult \citep{uciadult}     & 48\,842 & 14 & Income $>\!\$50\text{k}$ & Binary classification \\
Boston Housing \citep{harrison1978boston}\footnotemark & 506 & 13 & Median price   & Regression \\
California Housing \citep{pace1997california} & 20\,640 & 8 & Median value   & Regression \\
\bottomrule
\end{tabular}
\end{table}
\footnotetext{The Boston Housing dataset has known ethical
concerns \citep{caroll2020boston}; substitute California Housing or
another regression dataset if preferred.}

\section{Experimental Setup}
\label{sec:setup}

For every experiment the following invariants hold:

\begin{itemize}[leftmargin=1.4em]
  \item Random seed: \code{random\_state=42} throughout.
  \item Split: 80/20 train/test via
        \code{train\_test\_split} with stratification for
        classification when every class has $\ge 2$ samples.
  \item Cross-validation: $k=3$ folds inside
        \code{cross\_val\_score} for the bake-off; $k=3$ for the
        Optuna sweep.
  \item Candidates: RandomForest ($n\!=\!200$), XGBoost
        ($n\!=\!200$), LightGBM ($n\!=\!200$),
        LogisticRegression or Ridge as baseline. All defaults from
        \file{mcp\_servers/modeling.py::\_candidates}.
  \item Metrics: \code{f1\_weighted} for classification, $R^2$ for
        regression.
  \item Preprocessing: shared \code{ColumnTransformer}
        (median-impute + scale for numerics, most-frequent + one-hot
        for categoricals).
\end{itemize}

\section{Hardware and Software Environment}
\label{sec:env}

\begin{table}[H]
\centering
\caption{Reference execution environment used during development
smoke-testing. Users should record their own environment in
Appendix~\ref{app:extra} when they run the protocol.}
\label{tab:env}
\begin{tabularx}{\linewidth}{@{}lX@{}}
\toprule
\textbf{Item} & \textbf{Reference value} \\
\midrule
OS            & Windows 11 (development), Linux/macOS supported \\
Python        & 3.11+ \\
CPU           & Consumer laptop, 4--8 physical cores \\
RAM           & 16~GB \\
GPU           & Not required \\
LLM provider  & Groq Chat Completions API \\
LLM model     & \code{llama-3.3-70b-versatile} \\
Streamlit     & $\ge$ 1.36 \\
FastMCP       & $\ge$ 3.0 \\
scikit-learn  & Latest stable at install time \\
XGBoost / LightGBM & Latest stable at install time \\
\bottomrule
\end{tabularx}
\end{table}

\section{Baselines}
\label{sec:baselines}

The natural baselines for the platform are:

\begin{itemize}[leftmargin=1.4em]
  \item \textbf{Notebook baseline.} A manually written notebook that
        loads the dataset, applies the same preprocessing, trains
        each candidate with default parameters, and reports the
        same metrics. This baseline should produce identical
        numbers to \projname{} when the same random seed is used;
        any deviation indicates a bug.
  \item \textbf{Auto-sklearn} \citep{feurer2015autosklearn}.
        A stronger AutoML baseline. \projname{} is not expected to
        outperform Auto-sklearn on accuracy; the comparison is
        primarily on \emph{usability} and \emph{artefact quality}.
  \item \textbf{AutoGluon-tabular} \citep{erickson2020autogluon}.
        Another strong baseline that produces stacked ensembles.
\end{itemize}

\section{Metrics Reported}
\label{sec:metrics-reported}

For each dataset, the protocol records:

\begin{enumerate}[leftmargin=1.4em]
  \item Champion model name.
  \item Primary metric value on the test split.
  \item CV mean and standard deviation of the primary metric.
  \item Full leaderboard (all candidates with metric and training time).
  \item Wall-clock time for the bake-off tool call.
  \item Number and total size of generated artefacts.
\end{enumerate}

Chapter~\ref{ch:results} discusses the interpretation of these
numbers.

\section{Reproducibility Checklist}
\label{sec:repro-checklist}

Following \citet{pineau2021reproducibility}, each experimental run
should record:

\begin{itemize}[leftmargin=1.4em]
  \item Dataset URL and SHA-256 hash of the CSV.
  \item Exact tool invocations and arguments (available in
        \file{data/artifacts/\_index.json}).
  \item \code{pip freeze} output.
  \item Random seed (\code{random\_state=42}).
  \item Wall-clock start and end times.
\end{itemize}

The platform provides most of these automatically; the dataset hash
and \code{pip freeze} are the only manual steps.

\section{Experimental Template}
\label{sec:template}

Table~\ref{tab:experiment-template} is the template to fill in per
dataset. Cells marked ``---'' are to be populated by the person
running the protocol.

\begin{table}[H]
\centering
\caption{Experimental result template (\exper{}). Populate each
row after running the protocol on the corresponding dataset.}
\label{tab:experiment-template}
\small
\begin{tabularx}{\linewidth}{@{}lXcccc@{}}
\toprule
\textbf{Dataset} & \textbf{Champion} & \textbf{Primary} & \textbf{CV mean} & \textbf{CV std} & \textbf{Time (s)} \\
\midrule
Iris & --- & --- & --- & --- & --- \\
Wine & --- & --- & --- & --- & --- \\
Breast Cancer & --- & --- & --- & --- & --- \\
Adult & --- & --- & --- & --- & --- \\
California Housing & --- & --- & --- & --- & --- \\
\bottomrule
\end{tabularx}
\end{table}

\section{LLM Behavioural Evaluation}
\label{sec:llm-eval}

In addition to the numerical AutoML metrics, the platform should be
evaluated on \emph{agent behaviour}. Table~\ref{tab:llm-eval-rubric}
proposes a rubric.

\begin{table}[H]
\centering
\caption{Rubric for LLM behavioural evaluation. Score each item on
1--5 across a set of natural-language prompts.}
\label{tab:llm-eval-rubric}
\small
\begin{tabularx}{\linewidth}{@{}lX@{}}
\toprule
\textbf{Dimension} & \textbf{Description} \\
\midrule
Tool selection accuracy & Did the LLM pick the correct tool(s) for the request? \\
Argument fidelity & Are argument values (\code{file\_path},
\code{target\_column}, \code{model\_family}, \dots) correct? \\
Composition correctness & For multi-tool requests, did the sequence
match the pipeline (e.g.\ SHAP after bake-off)? \\
Result grounding & Are prose claims grounded in tool outputs, not
hallucinated? \\
Failure handling & When a tool errored, did the LLM recover
gracefully (retry, ask user, alternative tool)? \\
Formatting quality & Are outputs concise, well-structured, and free
of spurious markdown? \\
\bottomrule
\end{tabularx}
\end{table}

The rubric can be applied to a fixed suite of prompts (e.g.\ 20
prompts $\times$ 5 datasets) and averaged. This evaluation is
important because the platform's usability, not its raw accuracy, is
its primary contribution.
